\chapter{Conclusion}

This report described the first iteration of Granite, a self-hostable workbench for investigating silent failures in multi-agent LLM workflows. Starting from a research gap that the failure-attribution literature names itself, silent failures that leave no anomaly signal, and from the panel's feedback that our proposal described mechanisms rather than a product, we redefined the project around one workflow and one unit of value, the verified finding. In the first iteration we built and deployed a working slice: a five-agent refund workflow on public retail data, executable checks against a source of truth, deterministic provenance attribution, verification by controlled replay with controls under a fixed verdict rule, human review and tamper-evident evidence. The system is covered by 381 automated tests and continuous integration, its design artefacts are generated from the code, and a pre-registered feasibility run met every threshold while catching a genuine model error and exposing one design limit.

The results are deliberately modest in scope: one workflow, one fault type, one small model and ten tasks, with no baseline comparison or user study yet. Process gaps remain as well, chiefly visible per-member ownership of the code, which we are addressing through reviewed pull requests and code walkthroughs. Iteration 2 broadens the inputs, workflows and fault types, adds repair and recovery, and produces the first evaluation against baselines, so that by FYP-1 Final the claims in this report can be tested on a labelled benchmark with clean controls.
